\section{Resultados}
\label{sec:resultados}

A Figura~\ref{fig:antesdepois} mostra o antes e o depois nos dois vídeos
e na webcam. No painel da webcam aparece o texto que o programa escreve
na tela: o fundo atual e o número de objetos.

\begin{figure}[H]
\centering
\includegraphics[width=0.9\linewidth]{fig_antes_depois.jpg}
\caption{Antes e depois. Em cada grupo, da esquerda para a direita:
quadro original, máscara e resultado. (a) Vídeo claro, quadro 520; (b)
vídeo escuro, quadro 580, com o braço separado do corpo contado como
segundo objeto; (c) webcam, na demonstração gravada.}
\label{fig:antesdepois}
\end{figure}

A Tabela~\ref{tab:acertos} traz o resultado da avaliação nos 12 quadros
com gabarito.

\begin{table}[H]
\centering
\caption{IoU entre a máscara do programa e o gabarito pintado à mão.
Acerto quando IoU $\geq 0{,}8$.}
\label{tab:acertos}
\small\setstretch{1.1}
\begin{tabular}{llllc}
\toprule
Vídeo & Quadro & Situação & IoU & Resultado \\
\midrule
Claro  & 520 & entrando, em pé                  & 0,844 & acerto \\
Claro  & 600 & sentado, ventilador atrás         & 0,751 & erro \\
Claro  & 760 & acenando, ventilador atrás        & 0,753 & erro \\
Claro  & 820 & braço levantado                   & 0,787 & erro \\
Claro  & 880 & em pé, ventilador atrás           & 0,772 & erro \\
Claro  & 940 & bem perto da câmera               & 0,829 & acerto \\
\midrule
Escuro & 400 & entrando, borrado de movimento    & 0,748 & erro \\
Escuro & 460 & sentado                           & 0,783 & erro \\
Escuro & 540 & acenando                          & 0,741 & erro \\
Escuro & 580 & braço levantado                   & 0,816 & acerto \\
Escuro & 680 & em pé, perto                      & 0,791 & erro \\
Escuro & 720 & bem perto da câmera               & 0,875 & acerto \\
\midrule
\multicolumn{3}{l}{\textbf{Total: 4 acertos em 12 (33,3\%)}} & \textbf{0,791} & (média) \\
\bottomrule
\end{tabular}
\end{table}

O sistema acertou 2 dos 6 quadros em cada vídeo, \textbf{4 de 12 no
total (33,3\%)}, com IoU médio de \textbf{0,791}. A taxa baixa esconde
um resultado mais equilibrado: nenhum quadro ficou abaixo de 0,74, e
três erros ficaram entre 0,78 e 0,80, logo abaixo do critério.

O erro é quase todo de um tipo só (Figura~\ref{fig:avaliacao}): o
programa deixou de fora de 12\% a 26\% da pessoa em cada quadro, mas
quase nunca marcou fundo como pessoa (no máximo 2,6\% a mais). As
falhas se concentram na camiseta escura na frente do ventilador preto e
no short na frente da cadeira, a camuflagem prevista na
Seção~\ref{sec:mog2}.

\begin{figure}[H]
\centering
\includegraphics[width=0.75\linewidth]{fig_avaliacao.jpg}
\caption{Máscara do programa comparada com o gabarito. Verde: acerto;
vermelho: pessoa que o programa deixou de fora; azul: fundo marcado
como pessoa. Nos quadros claro 600 e 820, o buraco redondo é a camiseta
na frente do ventilador.}
\label{fig:avaliacao}
\end{figure}

A contagem de objetos acertou em \textbf{75 dos 78 quadros salvos
(96,2\%)}. Nos 36 quadros sem ninguém, ela deu 0; nos 42 com a pessoa,
deu 1 em 39. Os três erros (claro 820 e escuro 580 e 660) contaram 2
objetos, porque o braço levantado se separou do corpo na máscara.
